% from ega0-19b-fr.tex
\oldpage[0\textsubscript{IV-1}]{116}

\begin{examples}[19.10.3]
\label{0.19.10.3-fr}
Nous nous bornerons aux anneaux \emph{discrets}.
\begin{enumerate}
  \item[(i)] Si l'homomorphisme structural $\rho:A\to B$ est \emph{surjectif},
  $B$ est une $A$-algèbre \emph{formellement non ramifiée}, car avec les
  notations de (19.10.2), l'homomorphisme composé
  $A\xrightarrow{j}B\xrightarrow{v}E$ doit être l'homomorphisme structural,
  donc $v$ est unique (s'il existe). Par contre, $B$ n'est formellement lisse
  que si $\rho$ est \emph{bijectif}.
  \item[(ii)] Soient $A$ un anneau, $S$ une partie multiplicative de $A$~;
  alors $B=S^{-1}A$ est une $A$-algèbre \emph{formellement étale}. En effet,
  soient $j:A\to B$ l'homomorphisme canonique, $\rho:A\to E$ un homomorphisme
  d'anneaux, $\mathfrak J$ un idéal nilpotent de $E$,
  $p:E\to C=E/\mathfrak J$ l'homomorphisme canonique. Si $u:B\to C$ est un
  $A$-homomorphisme, $u(s/1)$ est inversible dans $C$ pour $s\in S$ puisque
  $s/1$ est inversible dans $B$~; mais $u(s/1)=p(\rho(s))$, et comme
  $\mathfrak J$ est nilpotent, le fait que l'image de $\rho(s)$ dans
  $E/\mathfrak J$ soit inversible entraîne que $\rho(s)$ lui-même est
  inversible dans $E$ ($\mathbf 0_{\mathrm I}$, 7.1.12)~; on en conclut que
  $\rho$ se factorise d'une seule manière en
  $A\xrightarrow{j}B\xrightarrow{v}E$, et l'on a évidemment
  $u\circ j=(p\circ v)\circ j$, d'où $u=p\circ v$, ce qui prouve notre
  assertion.
  \item[(iii)] Une extension finie séparable $k'$ d'un corps $k$ est une
  $k$-algèbre formellement lisse (19.6.1)~; nous verrons plus loin que pour
  qu'une extension de \emph{type fini} $K$ d'un corps $k$ soit une $k$-algèbre
  formellement étale, il faut et il suffit qu'elle soit formellement non
  ramifiée, et cela équivaut aussi au fait que $K$ est une extension finie et
  séparable de $k$ (21.7.4).
  \item[(iv)] Une algèbre de polynômes $A[X_\alpha]_{\alpha\in I}$ sur un
  anneau $A$ est \emph{formellement lisse} sur $A$ (19.3.3) mais \emph{non
  formellement étale} si $I$ n'est pas vide, comme il résulte aussitôt des
  définitions.
\end{enumerate}
\end{examples}

\begin{remark}[19.10.4]
\label{0.19.10.4-fr}
Le même raisonnement que dans (19.4.3) montre que pour vérifier qu'une
$A$-algèbre topologique $B$ est formellement non ramifiée, on peut, dans la
définition (19.10.2), se borner au cas où le noyau de $p:E\to C$ est de
\emph{carré nul}~; on peut en outre se borner au cas où $u(B)\subset p(E)$
(sans quoi il n'y a aucun $v$ factorisant $u$) et, en remplaçant $C$ par $u(B)$
et $E$ par $p^{-1}(u(B))$, supposer $u$ et $p$ surjectifs.
\end{remark}

% (continued in ega0-20-fr.tex, before its first page marker)
\section{Dérivations et différentielles}
\label{section:0.20-fr}
\label{0.20-fr}

Les notions introduites dans ce paragraphe seront reprises sous forme
géométrique au chapitre~IV, §~16, et joueront un rôle important dans l'étude
des préschémas. Leur importance dans ce chapitre tient d'abord à leurs
relations avec la notion d'algèbre formellement lisse, et notamment aux
théorèmes (20.4.9), (20.5.7) et (20.5.12), qui seront traduits en langage
géométrique dans le paragraphe du chapitre~IV consacré aux morphismes lisses,
et sont d'un usage constant dans les applications. D'autre part, les notions
différentielles serviront à démontrer au §~22 d'importants critères de
régularité qui joueront un rôle essentiel dans l'étude approfondie des anneaux
locaux noethériens faite au §~7 du chapitre~IV.

\subsection{Dérivations et extensions d'algèbres}
\label{subsection:0.20.1-fr}

